% !TeX program = xelatex
% ============================================================
% research/report/report.tex — Отчёт исследовательского курсового
% проекта ОП «Программная инженерия» ФКН НИУ ВШЭ.
%
% Оформление: ГОСТ 7.32-2017 «Отчёт о НИР». Чистый документ БЕЗ штампа
% (в отличие от ЕСПД-документов прикладного курсового). Структура и язык
% ветвятся по project.lang:
%   RU — форма ОП ПИ (ru pi): Список исполнителей → Реферат → Содержание →
%        Термины → Введение → главы (обзор→метод→результаты, «Выводы по
%        главе») → Заключение → Список источников (ГОСТ Р 7.0.5-2008) →
%        Приложения;
%   EN — англ. академическая форма ФКН: Annotation → Keywords → Contents →
%        Introduction → Literature review → Methodology → Experiments &
%        results → Conclusion → References → Appendices.
%
% Как пользоваться: пишите свой текст вместо жёлтых \hseFill{…}; данные о
% себе, руководителе, группе и годе задайте в настройках проекта.
% ============================================================
\documentclass[12pt,a4paper]{extreport}
\input{preamble}
\begin{document}
\input{title}


% ============================================================
%  ОТЧЁТ О НИР по ГОСТ 7.32-2017 (project.lang == ru)
% ============================================================
\chapter*{СПИСОК ИСПОЛНИТЕЛЕЙ}
\addcontentsline{toc}{chapter}{Список исполнителей}
% ГОСТ 7.32-2017 §5.2. У курсового отчёта исполнитель один — студент.
\noindent
Ответственный исполнитель, студент \hseCourse{} курса \hseFaculty,\\
образовательная программа «\hseSpecialization», группа \hseGroup

\vspace{0.4cm}
\noindent
\rule{5cm}{0.4pt}\quad\rule{3cm}{0.4pt}\hfill\hseAuthorName\\
{\scriptsize подпись, дата}

\clearpage
% ============================================================
\chapter*{РЕФЕРАТ}
\addcontentsline{toc}{chapter}{Реферат}
% ============================================================
% ГОСТ 7.32-2017 §5.3: сведения об объёме + ключевые слова + текст.
Отчёт \pageref{LastPage}~с., \ldots~рис., \ldots~табл., \ldots~источн., \ldots~прил.

\medskip
% Ключевые слова: 5–15 слов/словосочетаний, ПРОПИСНЫМИ, в именительном
% падеже, в строку через запятую, без точки в конце (§6.12.2).
\hseFill{КЛЮЧЕВОЕ СЛОВО ОДИН, КЛЮЧЕВОЕ СЛОВО ДВА, КЛЮЧЕВОЕ СЛОВО ТРИ}

\medskip
\noindent\textbf{Объект исследования} — \hseFill{что изучается в целом}.

\noindent\textbf{Предмет исследования} — \hseFill{конкретная сторона объекта}.

\noindent\textbf{Цель работы} — \hseFill{одна формулировка достигаемого научного результата}.

\noindent\textbf{Задачи:} \hseFill{провести обзор; предложить метод; реализовать и провести эксперимент; проанализировать результаты}.

\noindent\textbf{Результаты работы} — \hseFill{кратко: что получено (метод, модель, ПО, оценки, выводы) и область применения}.

\clearpage
% ── Содержание ──
\tableofcontents
\clearpage

% ============================================================
\chapter*{ТЕРМИНЫ И ОПРЕДЕЛЕНИЯ}
\addcontentsline{toc}{chapter}{Термины и определения}
% ============================================================
В~настоящем отчёте о~НИР применяют следующие термины с~соответствующими
определениями.

\begin{description}[leftmargin=0pt, labelindent=0pt, font=\bfseries, style=nextline]
  \item[\hseFill{Термин 1}] \hseFill{определение}.
  \item[\hseFill{Термин 2}] \hseFill{определение}.
\end{description}

\clearpage
% ============================================================
\chapter*{ВВЕДЕНИЕ}
\addcontentsline{toc}{chapter}{Введение}
% ============================================================
% Актуальность, цель, задачи, объект, предмет, краткий обзор, структура.
Обоснуйте актуальность: какая проблема решается, кого она затрагивает и
почему существующие подходы недостаточны~\cite{gost732}.

\textbf{Цель работы} — \hseFill{что именно достигается}.

\textbf{Задачи:}
\begin{enumerate}[label=\arabic*), leftmargin=1.25cm]
  \item \hseFill{провести обзор существующих подходов};
  \item \hseFill{предложить и формализовать метод (модель, алгоритм)};
  \item \hseFill{реализовать программное средство и провести эксперимент};
  \item \hseFill{проанализировать результаты и сделать выводы}.
\end{enumerate}

\textbf{Объект исследования} — \hseFill{что изучается в целом}.
\textbf{Предмет исследования} — \hseFill{конкретная сторона объекта}.

\clearpage
% ============================================================
\chapter{АНАЛИЗ ПРЕДМЕТНОЙ ОБЛАСТИ И ОБЗОР}
% ============================================================
\section{Современное состояние проблемы}
Дайте аналитический обзор предметной области и существующих подходов со
ссылками~\cite{gost732}.

\section{Сравнительный анализ существующих подходов}
Сопоставьте известные методы по значимым критериям; результат сравнения
приведён в~таблице~\ref{tab:approaches}.

\begin{table}[H]
\caption{Сравнение существующих подходов}
\label{tab:approaches}
\centering
\begin{tabular}{|>{\raggedright\arraybackslash}p{5cm}|c|c|c|}
\hline
\textbf{Критерий} & \textbf{Подход~А} & \textbf{Подход~Б} & \textbf{Предлагаемый} \\
\hline
Критерий~1 & \cmpplus  & \cmpminus  & \cmpplus \\
\hline
Критерий~2 & \cmpminus & \cmpplanned & \cmpplus \\
\hline
\end{tabular}
\end{table}

\section{Выводы по главе}
Сформулируйте нерешённую проблему и постановку задачи исследования.

\clearpage
% ============================================================
\chapter{МЕТОД (МОДЕЛЬ, АЛГОРИТМ) РЕШЕНИЯ}
% ============================================================
\section{Формальная постановка}
Введите обозначения и строго сформулируйте задачу (вход, выход, ограничения,
критерий качества, например целевую функцию $f(x)\to\min$).

\section{Предлагаемый метод}
Опишите предлагаемый метод (модель, алгоритм): основные шаги, допущения,
параметры.

\section{Выводы по главе}
Сформулируйте ключевые результаты главы.

\clearpage
% ============================================================
\chapter{ЭКСПЕРИМЕНТАЛЬНОЕ ИССЛЕДОВАНИЕ И РЕЗУЛЬТАТЫ}
% ============================================================
% Для инженерной работы без вычислительного эксперимента эту главу можно
% заменить на «РЕАЛИЗАЦИЯ И АПРОБАЦИЯ».
\section{Методика, данные и метрики}
Опишите план эксперимента, используемые данные (объём, источник, подготовка)
и метрики оценки качества.

\section{Результаты и их анализ}
Приведите результаты (таблица~\ref{tab:results}), сравните с базовыми методами
и оцените значимость различий.

\begin{table}[H]
\caption{Результаты эксперимента}
\label{tab:results}
\centering
\begin{tabular}{|>{\raggedright\arraybackslash}p{5cm}|c|c|}
\hline
\textbf{Метод} & \textbf{Метрика~1} & \textbf{Метрика~2} \\
\hline
Базовый метод      & \ldots & \ldots \\
\hline
Предлагаемый метод & \ldots & \ldots \\
\hline
\end{tabular}
\end{table}

\section{Выводы по главе}
Сформулируйте итоги экспериментального исследования.

\clearpage
% ============================================================
\chapter*{ЗАКЛЮЧЕНИЕ}
\addcontentsline{toc}{chapter}{Заключение}
% ============================================================
По результатам выполнения работы:
\begin{enumerate}[label=\arabic*), leftmargin=1.25cm]
  \item \hseFill{итог по задаче 1};
  \item \hseFill{итог по задаче 2}.
\end{enumerate}
Цель работы достигнута. Перспективы дальнейшей работы: \hseFill{направления}.

\clearpage
% ── Список источников по ГОСТ Р 7.0.5-2008 (методичка ФКН/ПИ): единый
%    алфавитный нумерованный список, сначала латиница, затем кириллица;
%    разделитель областей «. —»; сетевые ресурсы — «[Электронный ресурс].
%    — URL: … (дата обращения: ДД.ММ.ГГ)». ──
\phantomsection
\addcontentsline{toc}{chapter}{Список использованных источников}
\begin{thebibliography}{99}
  \bibitem{article}
  Фамилия, И.~О. Название статьи / И.~О.~Фамилия // Название журнала. --- 2024. --- № 1. --- С.~10--20.
  \bibitem{webresource}
  Название ресурса [Электронный ресурс]. --- URL: \url{https://example.org} (дата обращения: 01.01.2026).
  \bibitem{gost732}
  ГОСТ 7.32--2017. Отчёт о~научно-исследовательской работе. Структура и~правила оформления. --- М.: Стандартинформ, 2017. --- 27~с.
\end{thebibliography}

\clearpage
% ── Приложения (ГОСТ 7.32-2017 §6.14), рус. буквами А, Б, В… — опционально ──
\appendix
\renewcommand{\thechapter}{\Asbuk{chapter}}
\titleformat{\chapter}[hang]
  {\normalfont\large\bfseries\centering}{Приложение~\thechapter}{1ex}{}
\chapter{Иллюстрации}
Разместите здесь объёмные рисунки, листинги и таблицы данных.


\end{document}
